\itemhead{习题1.1 · 第6题（教材p.12）}{XJ-B1-C01-HW-P012-Y1286}
\label{XJ-B1-C01-HW-P012-Y1286}
{\footnotesize 湘教版教材p.12 / PDF第19页 / 习题1.1 / 第6题}
\field{原题}{\subq{1}只有一个元素的集合有多少个子集？
\subq{2}两个元素组成的集合有多少个子集？
\subq{3}$n\ (n\in\mathbb Z_+)$ 个元素组成的集合有多少个子集？}
\field{参考答案}{（1）2个；（2）4个；（3）$2^n$ 个。每个元素有选入或不选入两种独立选择。}
{\small \field{关联知识}{S03 子集、真子集与集合相等（教材帮p.11、12、13 / PDF20、21、22）}
\field{方法与易错点}{用逐元素检验或数轴包含；参数题先确认是否允许空集，再讨论端点。 易错：遗漏空集；把元素个数相等当作集合相等；把子集符号当作属于符号。}
\field{答案依据}{答案册相关题组：PDF第4、5页（册内页3、4）；按当前题库原题逐项重算，不按旧版题号机械对套。}}
